% !TEX encoding = UTF-8 Unicode
% !TEX root = ../Saunkin.tex

\chapter{Implementation and evaluation} \label{c4:implementation-evaluation}

\section{Implementation} \label{sec:implementation}

The system implements the citation analysis described in Chapter 3.
It takes an author name as input and returns citation counts, broken down into direct, co-author or external citations.
This section describes the component structure, data handling and processing steps to produce the results.

\subsection{System overview} \label{sec:system-overview}

The system is intended to run on a low-end machine.
It provides a web interface for the user (frontend) and a backend service, encapsulating the business logic, that collects and processes all data.
Both parts are written in \textsf{JavaScript}: the interface is built on \textsf{React}, and the backend targets \textsf{Node.js}.
The frontend sends requests and displays results, while the backend communicates with the publication services to perform citation analysis.

The first use case is to identify the author whose citations user wants to examine.
The search returns matching profiles, with affiliations where available, and lets the user inspect their publications before requesting metrics for a specific author.
The profiles are grouped by provider as each source uses identifiers not fully compatible with other sources.
Matching a profile from one source to a profile from another would require a separate identity matching process, which is out of scope for this work.

The second use case, having selected an author, is for the backend to collect publications and incoming citations.
It compares the author lists of each citing publication to classify citations and, finally, calculate totals.
The results, returned to the frontend, contain both the aggregate counts and the raw data.

Collecting these data can take time, especially for sources with low throughput and for authors with many publications.
The backend performs the retrieval as a background job, while the UI polls and displays its progress.
To avoid repeating the same work on every visit, the system keeps the collected author/publication/contribution/citation data, as well as caches provider responses in an \textsf{SQLite} database.
This keeps storage local, without a separate database server, and reduces the number of requests needed for later searches and calculations.

The architecture of the backend follows the Model-View-Controller (MVC) pattern.
Controllers receive user requests and pass them to services, which coordinate the collection and analysis.
Provider connectors handle the different external sources fetching and processing caveats.
Repositories read from and write to the database tables.
This separation keeps provider-specific details and storage operations apart from the citation classification rules.

The frontend serves the purpose of UI alone, and follows a different structure.
Components display individual parts of the interface, containers combine them and manage data requests, and pages arrange the views around navigation paths.
This gives the interface a maintainable structure while keeping citation analysis in the backend.

The following sections describe the provider choices, data model and processing steps in more detail.

\subsection{Providers and limitations} \label{sec:providers}

\subsection{Data} \label{sec:data}

\subsection{Storage and caching} \label{sec:storage-caching}

\subsection{Processing pipeline} \label{sec:processing-pipeline}

\subsection{User interface} \label{sec:user-interface}

\subsection{Testing} \label{sec:testing}

\section{Evaluation} \label{sec:evaluation}

\subsection{Citation classification} \label{sec:classification-checks}

\subsection{Results for selected researchers} \label{sec:author-case-studies}

\subsection{Comparison of data sources} \label{sec:data-source-comparison}
